% New result after the fixed-depth certification; independent review pending.
\subsection{Finite first failures are unbounded}

\begin{result}{C.2.13}{Proved}{Selected fixed-depth ghosts have a finite later failure}
\label{thm:ghost-unbounded-finite-failures}
For every \(R\ge1\) and positive integer \(N\),
the construction of Theorem~\ref{thm:ghost-fixed-depth-survival}
can be made to give a pair \((A_{R,N},B_{R,N})\)
whose first empty unrestricted integral nilpotent
section fibre occurs at a finite class strictly
larger than \(R\). It has no section of the full
geometrically pro-five arithmetic fundamental group.
The finite excluding class is asserted to exist;
no formula for its least value is asserted.

One fixed finite etale cover detects the full
pro-five exclusion of all these selected examples:
the point-three twist of the characteristic
Heisenberg cover of geometric degree \(125\).
The exact cover subgroup and the local comparison
are specified below. The statement concerns these
selected outputs of A.65, rather than every output
of that construction.
\end{result}
\begin{proof}
Replace \(R\) by \(\max(R,3)\) in the construction.
Choose its integer \(C\) also divisible by five.
Every fresh exponent in (VX.6) is then divisible
by five. We impose one additional compatible
condition on its first positive support prime.

Use the tangentially marked class-two exponent-five
quotient from the proof of
Theorem~\ref{thm:ghost-first-integral-class3}:
\[
 E_{2,5}=H_5\rtimes_{\chi}G_{\mathbb Q},
 \qquad H_5=N_{2,5}/5,
 \tag{UF.1}
\]
where \(H_5\) has order \(125\), abelianization
\(\mathbb F_5(1)^2\), and center
\(\mathbb F_5(2)\). The action in (UF.1) is
the pure cyclotomic class-two action in the
literal Hall coordinates of
Corollary~\ref{cor:ghost-class2-lift}.
Its finite arithmetic image is
\(H_5\rtimes\mathbb F_5^\times\).
Let \(L\) be the specialization at three of
the marked generic cover
\[
 X^5=t,\quad Y^5=1-t,\quad
 Z^5=\Theta(X)=\prod_{i=0}^4(1-\zeta^iX)^i,
 \tag{UF.2}
\]
over \(\mathbb Q(\zeta)\), \(\zeta^5=1\),
\(\zeta\ne1\), with the tangential branch
\(Y(0)=Z(0)=1\). The exact identification of
(UF.2) with this marked quotient, including
its descent, was proved in that theorem.
In particular \(\operatorname{Gal}(L/\mathbb Q)
 =H_5\rtimes\mathbb F_5^\times\), its maximal
abelian subfield is \(\mathbb Q(\mu_5)\), and
the maximal abelian subextension over
\(\mathbb Q(\mu_5)\) is obtained by adjoining
\(3^{1/5},(-2)^{1/5}\).

At the first arithmetic stage put \(n=J!\),
and let \(M\) be its full triangular CRT
modulus. Its selected unit class is one
modulo \(n\). In particular it is one
modulo five, although its prescribed action
on \(\mathbb Q(\mu_M)\) need not be the identity.
Set
\[
 F=\mathbb Q(\mu_M),\qquad
 R_b=F((-2)^{1/n}),\qquad
 R'=R_b(3^{1/5}).
 \tag{UF.3}
\]
The original CRT automorphism \(\sigma_F\)
extends to \(R_b\) while fixing \((-2)^{1/n}\),
by the exact intersection proved in
Theorem~\ref{thm:ghost-general-exponent-arithmetic}.
The fifth classes of \(3,-2\) are independent
in \(F\). Indeed a nontrivial rational
combination has a valuation not divisible
by five, and its pure fifth-root field has
degree five and is nonnormal over
\(\mathbb Q\); it cannot lie in a cyclotomic
field. The cyclic extension \(R_b/F\) has
unique degree-five subfield
\(F((-2)^{1/5})\), so
\(R_b\cap F(3^{1/5})=F\). Thus
\(\sigma_F\) extends to \(\sigma_{R'}\)
fixing both radicals in (UF.3).

Since \(L\cap F=\mathbb Q(\mu_5)\),
\(\operatorname{Gal}(LF/F)=H_5\).
The extension \(R'/F\) is abelian and
contains its full abelian subextension.
Consequently the exact intersection is
\[
 LF\cap R'=F(3^{1/5},(-2)^{1/5}).
 \tag{UF.4}
\]
Choose a nonzero central element \(z\in H_5\).
The automorphisms \(z\) on \(L\) and
\(\sigma_F\) on \(F\) agree on
\(\mathbb Q(\mu_5)\). Their resulting
automorphism on \(LF\) agrees with
\(\sigma_{R'}\) on (UF.4): both fix its
two displayed radicals and act by
\(\sigma_F\) on \(F\).
This proves the required fibre-product
compatibility. Chebotarev supplies a
first prime \(q\) meeting every original
CRT and splitting requirement, also
splitting \(3^{1/5}\), whose point-three
Frobenius is a nonzero central element
of \(H_5\). Its conjugates may multiply
that element by a nonzero scalar; they
remain nonzero central elements. Continue
the arithmetic recursion unchanged for
\(r,s\) and all subsequent stages.

\emph{All integral class-two lifts have the same finite image.}
Choose literal abelian cocycles for the
products (VX.7), so their reductions modulo
five equal those of \((3,-2)\). The
central difference from the point-three
class-two section is then a finite
\(\mathbb F_5(2)\)-cocycle. Its inertia
vanishes at every fresh support prime:
the positive patterns give a multiple
of their exponent, while the negative
pattern gives a binomial expression
and an integral multiple of that
exponent. Both vanish modulo five.
At every zero-valuation prime the
integral tame relation gives zero
inertia already integrally. At two
and three, the central difference has
inertia killed by \(q-1\), a unit
modulo five, so again it is zero.
It is thus unramified outside five.
The exact vanishing
\(H^1(G_{\mathbb Q,\{5,\infty\}},
       \mathbb F_5(2))=0\), proved in
Lemma~\ref{lem:ghost-prefix-five},
makes this difference a central
coboundary. Every integral class-two
lift therefore has the point-three
image in (UF.1), up to geometric
conjugacy. At the selected prime
\(q\equiv1\pmod5\), the central
coboundary is zero on Frobenius,
so its nonzero central value is
unchanged even in these literal
coordinates.

\emph{The local ordinary point in the same marking.}
Let \(t_q\) be the ordinary local point
of the constructed pair. Its valuation
\(E_0(q)>0\) is divisible by five.
The equality \(A/3\in(\widehat{\mathbb Q^\times})^5\)
and the extra splitting condition give
\(t_q\in(\mathbb Q_q^\times)^5\).
Choose \(X\in\mathbb Q_q\) with
\(X^5=t_q\); then \(v_q(X)>0\).
Here \(\mu_5\subset\mathbb Q_q\), and
every \(1-\zeta^iX\) is a principal
\(q\)-unit. So are \(1-t_q\) and
\(\Theta(X)\). Raising a principal
\(q\)-unit to its fifth power is an
automorphism, since \(q\ne5\).
Thus \(Y,Z\) in (UF.2) belong to
\(\mathbb Q_q\), and the entire
\(H_5\)-torsor over \(t_q\) splits.
Because the arithmetic action in
(UF.1) is trivial on \(G_{\mathbb Q_q}\),
its local point section is the identity
homomorphism \(G_{\mathbb Q_q}\to H_5\),
up to geometric conjugation, which
does not change that identity. This
uses the same tangential marking as
the nonzero point-three Frobenius;
no change of arithmetic splitting
has been made in the comparison.

A full geometrically pro-five section
above the pair would, at this nonzero
valuation prime, be geometrically
conjugate to the ordinary point
section by
Proposition~\ref{prop:ghost-local-boundary}.
Its image under the actual quotient
(UF.1) would consequently be trivial
there. The preceding central Frobenius
calculation contradicts this. Thus
the full pro-five fibre is empty.

The compact-fibre proof of
Theorem~\ref{thm:ghost-fixed-pair-compatible},
applied at the single coefficient
prime five, now implies that some
finite integral nilpotency class
already has empty pro-five fibre.
The full integral fibre at that
class is therefore empty as well.
Theorem~\ref{thm:ghost-fixed-depth-survival}
provides a full integral class-\(R\)
section for this same pair, so the
least empty full integral class
exists and is strictly larger than
the originally prescribed \(R\).

Finally define the promised cover.
Let \(K\) be the geometric kernel
of \(\Pi\to H_5\), and let \(Y_5\)
correspond to the inverse image of
\(s_3(G_{\mathbb Q})\) in
\(H_5\rtimes G_{\mathbb Q}\).
This open subgroup has geometric
index \(125\), has the rational
point above three, and its induced
geometric pro-five topology is full
because the index is a power of five.
Every hypothesized full pro-five
section above the constructed pair
can be geometrically conjugated into
this subgroup, by the just-proved
class-two comparison. Locally at
the selected \(q\), the fibre of
\(Y_5\to U\) is the actual coset
set \(H_5\). The point section's
identity homomorphism and the
nonzero central point-three
Frobenius give a nontrivial central
translation on this set, with no
fixed coset. This proves the
specified cover exclusion with
its actual section-to-subgroup map.
\end{proof}
